\begin{table}[h]
\centering
\begin{tabular}{|p{0.35\textwidth}|p{0.55\textwidth}|}
\hline
\textbf{Variables} & \textbf{Definition} \\
\hline

\texttt{Number of rooms} & Self-explanatory \\
\hline

\texttt{Average room size} & Average room size \\
\hline

\texttt{Floor} & Number of floors \\
\hline

\texttt{Age} & 
\begin{tabular}[t]{l}
Age $< 3$ \\
After the Tohoku earthquake (from 2011) \\
After the Kobe earthquake (1995--2010) \\
After the para-seismic law (1981--1994) \\
After the City Planning Law (1968--1980) \\
After the WW2 (1945--1967)
\end{tabular}
\\
\hline

\texttt{Building structure} &
\begin{tabular}[t]{l}
Wood \\
Concrete light \\
Concrete intermediate \\
Concrete solid \\
Concrete very solid
\end{tabular}
\\
\hline

\texttt{Distance to station} & Distance to the nearest station. In the Japanese case, stations are a good proxy for city centers. \\
\hline

\texttt{Density} & Density Inhabited Districts (DID) is defined by the Population Census, where areas with a population density of 4,000 people per square kilometer or higher are adjacent to each other and have a total population of 5,000 or more. Density takes the value 0 if the municipality is not defined as DID, and takes $\log(\text{density})$ otherwise. \\
\hline

\texttt{Net migration rate} & Number of in-migrants over number of total population in 2015. \\
\hline

\end{tabular}
\end{table}